\subsection{Directional Discourse Supervision for Knowledge-Map Generation}

Одним из наиболее содержательных исследовательских выводов, вытекающих из проделанной работы, является несоответствие между природой целевой задачи и формой части используемого supervision-сигнала. В проекте CMGN требуется построить ориентированный граф управляющих связей между предложениями документа. Формально модель предсказывает матрицу
\[
\hat{Y} \in [0,1]^{N \times N},
\]
где элемент $\hat{y}_{ij}$ характеризует степень того, насколько предложение $s_i$ структурно управляет предложением $s_j$. Следовательно, сама постановка задачи является направленной: отношение между $s_i$ и $s_j$ в общем случае несимметрично.

При этом один из использованных в проекте источников автоматической разметки --- матрица попарного косинусного сходства предложений, построенная на основе DistilBERT-представлений. Такой сигнал по своей природе отражает преимущественно семантическую близость, а не иерархическую или риторическую направленность. В результате supervision вида
\[
L^{\mathrm{bert}}_{ij} = \frac{\cos(h_i,h_j)+1}{2}
\]
задает мягкое симметричное сходство между предложениями, но не кодирует различие между отношениями типа ``источник $\rightarrow$ уточнение'' и ``уточнение $\rightarrow$ источник''. Для задачи построения ориентированной карты знаний это является принципиальным ограничением.

Экспериментальные наблюдения подтверждают данную гипотезу. В отчете было показано, что DistilBERT-разметка дает худшие результаты по сравнению с направленными источниками supervision, тогда как риторическая разметка, особенно в ограниченном режиме \texttt{RST only NS}, демонстрирует более устойчивое качество. Это указывает на то, что для построения иерархической карты знаний важен не столько сам факт семантической близости предложений, сколько наличие асимметричного структурного сигнала, отражающего отношение ядра и зависимого фрагмента.

На этой основе можно сформулировать исследовательскую гипотезу о \emph{направленном дискурсивном supervision}. Пусть для документа строятся несколько вспомогательных графов:
\[
A^{\mathrm{rst}} \in \{0,1\}^{N \times N}, \qquad
A^{\mathrm{coref}} \in \{0,1\}^{N \times N}, \qquad
A^{\mathrm{sem}} \in [0,1]^{N \times N},
\]
где $A^{\mathrm{rst}}$ задает directed-связи, полученные из риторической структуры текста, $A^{\mathrm{coref}}$ --- directed-кореферентные переходы между предложениями, а $A^{\mathrm{sem}}$ --- мягкий семантический prior.

Тогда совокупный teacher-graph может быть задан в виде
\[
A^{\mathrm{teach}} = \alpha A^{\mathrm{rst}} + \beta A^{\mathrm{coref}} + \gamma A^{\mathrm{sem}},
\]
где коэффициенты удовлетворяют условию
\[
\alpha > \beta > \gamma \ge 0,
\]
что отражает более высокую значимость направленного дискурсивного сигнала по сравнению с симметричной семантической близостью.

Вместо покомпонентной регрессии только к матрице labels можно использовать комбинированную целевую функцию
\[
L = \lambda_1 L_{\mathrm{dir}} + \lambda_2 L_{\mathrm{rank}} + \lambda_3 L_{\mathrm{cons}},
\]
где $L_{\mathrm{dir}}$ --- directed-loss для восстановления ориентированных ребер, $L_{\mathrm{rank}}$ --- ranking-loss, обеспечивающая более высокий score для правильного родительского предложения по сравнению с альтернативами, а $L_{\mathrm{cons}}$ --- регуляризатор согласованности между несколькими структурными представлениями.

Например, directed-компонента может быть задана как взвешенная бинарная кросс-энтропия:
\[
L_{\mathrm{dir}}
=
-\sum_{i \ne j}
w_{ij}
\left(
A^{\mathrm{teach}}_{ij}\log \hat{y}_{ij}
+
(1-A^{\mathrm{teach}}_{ij})\log(1-\hat{y}_{ij})
\right),
\]
где веса $w_{ij}$ могут усиливать вклад риторически наиболее значимых переходов.

Ранжирующая компонента может быть записана в виде
\[
L_{\mathrm{rank}}
=
\sum_{j=1}^{N}
\sum_{n \ne p_j^\ast}
\max\bigl(0,\; m - \hat{y}_{p_j^\ast j} + \hat{y}_{nj}\bigr),
\]
где $p_j^\ast$ --- правильный родитель узла $j$, а $m>0$ --- margin. Такая формулировка особенно уместна в задаче построения карты знаний, поскольку модель должна не просто повышать значения на ``хороших'' ребрах, а выбирать структурно предпочтительного родителя для каждого предложения.

Наконец, компоненту согласованности можно ввести как штраф за расхождение между различными структурными источниками:
\[
L_{\mathrm{cons}}
=
\sum_{i,j}
\left|
\hat{y}^{\mathrm{rst}}_{ij} - \hat{y}^{\mathrm{coref}}_{ij}
\right|.
\]

Научная новизна такого подхода состоит в переходе от симметричного similarity-based supervision к \emph{асимметричному структурному supervision}, согласованному с природой самой задачи. В этой постановке модель обучается восстанавливать не просто семантическую близость предложений, а directed-отношения типа ``ядро $\rightarrow$ развитие'', ``первое упоминание $\rightarrow$ повторное упоминание'', ``структурный центр $\rightarrow$ зависимый фрагмент''. Тем самым supervision становится ближе к содержательной цели генерации карты знаний.

\subsection{TTED-Aware Training for Hierarchical Knowledge Maps}

Второе перспективное направление исследований связано с несоответствием между функцией потерь, используемой при обучении модели, и метрикой, по которой затем оценивается итоговое качество построенной карты знаний. В текущей реализации базовая supervised-компонента задается как среднеквадратичное отклонение между предсказанной и целевой матрицами связей:
\[
L_{\mathrm{sup}}
=
\frac{1}{N^2}
\sum_{i=1}^{N}\sum_{j=1}^{N}
(\hat{y}_{ij} - y_{ij})^2.
\]
Такое обучение является по существу edge-wise: модель оптимизирует локальное соответствие отдельных элементов матрицы.

Однако итоговая оценка в проекте существенно богаче. Помимо sentence-level и keyword-level показателей используется TTED --- метрика, сравнивающая generated и gold mind-map структуры как деревья. В отчете подчеркивается, что TTED штрафует не только локальное несовпадение ребер, но и более сложные расхождения:
\begin{itemize}
    \item структурные ошибки в родительско-дочерних отношениях;
    \item семантическое несоответствие текстов сопоставляемых узлов;
    \item иерархическое несоответствие, когда узел оказывается на неверном уровне дерева.
\end{itemize}

Следовательно, возникает принципиальный \emph{train--evaluation mismatch}: модель обучается как предсказатель матрицы попарных связей, а оценивается как генератор иерархической смысловой структуры. Из этого естественно вытекает гипотеза о необходимости \emph{TTED-aware training}, то есть обучения, в котором функция потерь приближена к тем типам ошибок, которые действительно фиксирует TTED.

Наиболее реалистичной формой такого подхода является введение структурного суррогата, который не требует прямой дифференциации TTED, но учитывает его содержательный смысл. Пусть для каждого узла $j$ модель задает распределение вероятностей по возможным родителям:
\[
p(i \mid j)
=
\frac{\exp(z_{ij})}{\sum_{k \ne j}\exp(z_{kj})},
\]
где $z_{ij}$ --- score связи $i \rightarrow j$ до нормировки.

Если $p_j^\ast$ обозначает правильного родителя узла $j$ в gold-дереве, то parent-level компоненту можно определить как
\[
L_{\mathrm{parent}}
=
-\sum_{j=1}^{N}
\log p(p_j^\ast \mid j).
\]
В отличие от покомпонентной MSE, такая функция потерь непосредственно оптимизирует корректность выбора родителя, а значит, лучше согласуется с деревообразной природой итогового объекта.

Поскольку TTED чувствителен и к иерархическому положению узлов, разумно дополнительно учитывать глубину узла в дереве. Пусть $d_j^\ast$ --- глубина узла $j$ в gold-структуре, а $\hat{d}_j$ --- предсказание модели. Тогда можно ввести depth-loss:
\[
L_{\mathrm{depth}}
=
\sum_{j=1}^{N}
(\hat{d}_j - d_j^\ast)^2.
\]
Альтернативно глубина может рассматриваться как дискретная метка, и тогда применяется классификационная форма соответствующей функции потерь.

Еще одна важная составляющая TTED связана с сохранением правильной предковой структуры. Даже если локальный родитель выбран почти корректно, ошибка в верхней части дерева может привести к изменению целого пути узла к корню. Для учета этого эффекта можно задать бинарную матрицу предковых отношений
\[
A^\ast \in \{0,1\}^{N \times N},
\]
где
\[
A^\ast_{aj} =
\begin{cases}
1, & \text{если узел } a \text{ является предком узла } j,\\
0, & \text{иначе.}
\end{cases}
\]
Тогда ancestor-consistency loss принимает вид
\[
L_{\mathrm{anc}}
=
-\sum_{a,j}
\left[
A^\ast_{aj}\log \hat{A}_{aj}
+
(1-A^\ast_{aj})\log(1-\hat{A}_{aj})
\right],
\]
где $\hat{A}_{aj}$ --- предсказанная моделью вероятность предкового отношения.

Итоговая TTED-aware функция потерь может быть задана как
\[
L_{\mathrm{TTED}}
=
\lambda_1 L_{\mathrm{edge}}
+
\lambda_2 L_{\mathrm{parent}}
+
\lambda_3 L_{\mathrm{depth}}
+
\lambda_4 L_{\mathrm{anc}},
\]
где $L_{\mathrm{edge}}$ может совпадать со стандартной edge-level loss, а остальные компоненты приближают структурные аспекты TTED.

Содержательный смысл данной конструкции состоит в следующем. Обычная edge-level оптимизация отвечает на вопрос: \emph{верно ли предсказано отдельное ребро?} Напротив, TTED-aware обучение отвечает на более сильный вопрос: \emph{верно ли восстановлена иерархическая роль узла в дереве?} Именно это делает такой подход особенно уместным для генерации карт знаний, где качество определяется не только наличием связей, но и тем, как эти связи организованы в глобальную структуру.

Более сильная, но и более сложная версия данного направления состоит в переходе к minimum-risk training с reward, основанным непосредственно на TTED. Если $\hat{T}$ --- дерево, сгенерированное моделью, а $T^\ast$ --- gold-дерево, то можно задать риск
\[
L_{\mathrm{risk}}
=
\mathbb{E}_{\hat{T}\sim p_\theta(T\mid X)}
\bigl[
\mathrm{TTED}(\hat{T},T^\ast)
\bigr].
\]
На практике эта величина может аппроксимироваться по выборке candidate trees:
\[
L_{\mathrm{risk}}
\approx
\sum_{k=1}^{K}
p_\theta(\hat{T}_k\mid X)\,
\mathrm{TTED}(\hat{T}_k,T^\ast).
\]
Такая постановка концептуально особенно хорошо согласуется с архитектурой CMGN, поскольку в проекте уже присутствует RL-компонента в общей функции потерь. Следовательно, TTED может быть интерпретирован как более содержательный reward, ориентированный не на внешнее текстовое сходство, а на правильность итоговой структуры карты знаний.

С научной точки зрения основное достоинство TTED-aware training заключается в том, что оно устраняет разрыв между тем, чему модель обучается, и тем, по чему она затем оценивается. Если стандартная MSE учит модель приближать элементы матрицы попарных связей, то TTED-aware подход учит ее восстанавливать глобально согласованную иерархическую организацию документа. Именно это делает данное направление естественным и сильным кандидатом для последующей научной публикации.